\subsection{I-25. Intrinsic Grassmannian compression succeeds for invertible chunks}
\label{I-25}

The nested-image method is suggested by Liu (2026), Section~3.3,
printed pp.~7--8.  The following restricted refinement uses
the projective Grassmannian of one decomposable exterior row
as its ambient variety and is proved directly.

\paragraph{Statement.}
Let $A_i(Y_i)$ be $w\times w$ matrices over a field $L$, in
distinct variables, with every $A_i(0)$ invertible.  Fix a
nonzero decomposable row $p\in\bigwedge^r L^w$ and a linear
functional $\ell$ on $\bigwedge^rL^w$, where $1\le r<w$.
If
\[
 q(\mathbf Y)=\ell\bigl(p(\textstyle\bigwedge^r A_1(Y_1))
                  \cdots(\textstyle\bigwedge^r A_m(Y_m))\bigr)
 \ne0,
\]
then there is a set $S\subseteq[m]$ with
$|S|\le r(w-r)$ for which setting $Y_i=0$ outside $S$
keeps $q$ nonzero.  In particular $q$ has a nonzero monomial
of support at most $r(w-r)$, and $\operatorname{SV}_{r(w-r)}$
detects it.

For an $r$-minor test with an invertible-constant interior
chunk and at most one factor on either end, the whole minor
therefore has some nonzero monomial of support at most
$r(w-r)+2$.  This replaces the arbitrary compound-space bound
$\binom wr-1+2$ for this restricted class.  It bounds the
existence of a low-support monomial, not every monomial in the
lowest homogeneous component.

\paragraph{Proof.}
Use cut-dependent invertible constant changes of basis to make
every $A_i(0)=I$.  These changes arise from $\operatorname{GL}_w$
and hence preserve decomposability after taking exterior powers;
they adjust only the fixed endpoints.  Extend scalars to an
algebraic closure.  This does not change whether a polynomial
or its coordinate restriction is zero.

Restrict each parameter to the nonempty open set where its
factor is invertible.  Let $Z_i$ be the Zariski closure in
$\operatorname{Gr}(r,w)$ of the projective prefix states
\[
 [p(\textstyle\bigwedge^r A_1(Y_1))\cdots
        (\textstyle\bigwedge^r A_i(Y_i))].
\]
Each $Z_i$ is irreducible, as the closure of the image of an
irreducible product of open affine lines.  The value $Y_i=0$
is allowed and acts as the identity, so
$Z_{i-1}\subseteq Z_i$.  The starting variety $Z_0$ is one
point.  The Grassmannian has dimension $r(w-r)$: its dense
matrix chart consists of row spaces of $[I_r\ X]$ with
$X\in L^{r\times(w-r)}$.
Thus at most $r(w-r)$ of these inclusions are strict.

Whenever $Z_i=Z_{i-1}$, replace $A_i(Y_i)$ by $I$.
This preserves the closure of the prefix image and hence the
closure of every later image.  To justify the latter assertion,
on invertible factor values the exterior action is a regular
map on the projective Grassmannian; the closure of the image
under the next factor of a dense prefix set is the same as
the closure obtained from the whole prefix variety.  This
persists on composing successive factors.  Inductively one
may therefore remove every non-growth position while retaining
the final projective image closure.  At most $r(w-r)$ positions
remain.

The scalar $q$ vanishes identically exactly when that final
projective image lies in the hyperplane cut out by $\ell$.
It suffices to test this on the invertible-value open set,
since a polynomial vanishing there vanishes identically.
Preserving the image closure therefore preserves nonvanishing
of $q$.  A nonzero restriction to the retained coordinates
contains a monomial supported there, proving the first claim.

For the endpoint version, adjoin the at most two endpoint
variables to the coefficient field.  A selected row of an
$r$th compound endpoint is the wedge of $r$ rows of the
original endpoint matrix, hence is decomposable; it is nonzero
if the selected full minor is nonzero.  The right endpoint
supplies a linear functional.  The preceding argument applies
to the interior chunk over this rational function field.
The resulting restricted minor is still a polynomial in the
original variables and is nonzero there.  Restoring the at
most two endpoint positions gives the asserted support bound.
\hfill$\square$

\paragraph{Exact limitation.}
This construction exploits the exterior geometry instead of
treating $\bigwedge^r L^w$ as an arbitrary vector space of
dimension $\binom wr$.  It is an existence-and-enumeration
compression: the retained positions depend on the matrices,
and enumerating them or using the stated SV map still has an
exponent depending on $w$.  The proof also requires invertible
constant interior factors.  I-24 shows that replacing this
condition merely by preservation of all interval ranks at
least two cannot yield a polynomial-in-$w$ support bound.
The next useful rank-descent construction must accommodate
those singular resets and the original-variable cancellation
relations simultaneously; intrinsic dimension alone does
not supply that extension.
